\section*{ID9085: Relation: Dr 2)|r 1-R 2|ˆ3Lk=4π1C1C2r1r23(r1r2)(dr1×}
\paragraph{Metadata.}
PiID: 9658253375438 | RTEID: RTE:P35279.0726:T1.5804 | UUID: 55909bb9-8c08-5d2d-8d0f-e21083864a0a

\paragraph{Canonical Statement.}
lity; helps classify turbulence spectra. Falsifiability / Test Handle: Observed dependence of the dimensionless helicity ratio on system size would challenge this constraint. 0.0.45 ID 043: Helicity Linking Number Universality Type: Definition Formal Statement: The linking number LkLkLk between two magnetic flux tubes in a Trawinistic manifold is given by \$Lk=14πC1C2(r1r2)(dr1×\$ dr2)r1r23Lk = \textbackslash{}frac\{1\}\{4\textbackslash{}pi\}\textbackslash{}oint\_\{C\_1\}\textbackslash{}oint\_\{C\_2\}\textbackslash{}frac\{(\textbackslash{}mathbf\{r\} 1-\textbackslash{}mathbf\{r\} 2)\textbackslash{}cdot(\textbackslash{}text\{d\}\textbackslash{}mathbf\{r\} 1\textbackslash{}times \$\textbackslash{}text\{d\}\textbackslash{}mathbf\{r\} 2)\}\{|\textbackslash{}mathbf\{r\} 1-\textbackslash{}mathbf\{r\} 2|ˆ3\}Lk=4π1C1C2r1r23(r1r2)(dr1×\$ dr2) and is restricted to integer values. This integer is independent of system scale and boundary conditions. Interpretive Meaning: The linking number between magnetic structures is universally quantised. This universality underlies the robustness of helicity and topological constraints across scales. Dependencies: ID 041 Derivable Consequences: Predicts conservation of linking during reconnecti

\paragraph{Canonical Equation.}
\[
\text{d}\mathbf{r} 2)}{|\mathbf{r} 1-\mathbf{r} 2|ˆ3}Lk=4π1C1C2r1r23(r1r2)(dr1×
\]

\paragraph{Dictionary.}
Relation: Dr 2)|r 1-R 2|ˆ3Lk=4π1C1C2r1r23(r1r2)(dr1× — dr 2)|r 1-r 2|ˆ3Lk=4π1C1C2r1r23(r1r2)(dr1×

\paragraph{Encyclopedia.}
Relation: Dr 2)|r 1-R 2|ˆ3Lk=4π1C1C2r1r23(r1r2)(dr1× is an algebraic definition, written dr 2)|r 1-r 2|ˆ3Lk=4π1C1C2r1r23(r1r2)(dr1×. It is an algebraic definition expressing the quantity directly in terms of the variables on its right-hand side. It is built from r, L, k, C. Within the Trawin Zero Point registry it serves as one element of the framework's mathematical narrative.
